\documentclass[11pt]{article}

\usepackage[utf8]{inputenc}
\usepackage[T1]{fontenc}
\usepackage{geometry}
\usepackage{amsmath}
\usepackage{amssymb}
\usepackage{booktabs}
\usepackage{hyperref}
\usepackage{xurl}
\usepackage{xcolor}
\usepackage{listings}
\usepackage{enumitem}
\usepackage{parskip}
\usepackage{graphicx}
\graphicspath{{../images/}}

\geometry{margin=1in}

\lstset{
  language=Java,
  basicstyle=\ttfamily\small,
  keywordstyle=\color{blue},
  commentstyle=\color{gray},
  stringstyle=\color{teal},
  showstringspaces=false,
  columns=fullflexible,
  frame=single,
  framerule=0.2pt,
  breaklines=true
}

\setlist[itemize]{topsep=0.3em,itemsep=0.2em}
\setlist[enumerate]{topsep=0.3em,itemsep=0.2em}

% Simple per-section source line.
\newcommand{\notesources}[1]{\par\noindent{\small\color{gray}\textit{Sources:} \sloppy #1\par}}

% Unit-wise source strings for this subject (used by slides/notes).
% Path is relative to the lecture `latex/` folder (the compilation CWD).
% OOPS with Java (CSEG1044) — unit-wise sources
%
% These are short, student-facing reference strings intended to be used with:
%   - \slidesources{...}  (in beamer slides)
%   - \notesources{...}   (in lecture notes)
%
% Style inspiration: other subjects in My-Subjects/ use semicolon-separated sources
% and include an "accessed" date for web documentation.

\newcommand{\OOPSUnitISources}{Schildt, \textit{Java: A Beginner's Guide} (10e), McGraw-Hill (Java fundamentals + OOP basics).; Oracle Java Tutorials: Getting Started + Language Basics + Classes and Objects (accessed \today).; Java SE API docs (e.g., \texttt{String}, \texttt{Scanner}) (accessed \today).}

\newcommand{\OOPSUnitIISources}{Schildt, \textit{Java: The Complete Reference} (12e), McGraw-Hill (classes, inheritance, packages, interfaces).; Bloch, \textit{Effective Java} (3e), Addison-Wesley, 2018 (design choices; interfaces vs abstract classes).; Oracle Java Tutorials: Classes and Objects + Interfaces + Packages (accessed \today).; Java SE API docs (e.g., \texttt{Comparable}, \texttt{Runnable}) (accessed \today).}

\newcommand{\OOPSUnitIIISources}{Schildt, \textit{Java: The Complete Reference} (12e) (nested/inner classes, exceptions, threads, I/O).; Oracle Java Tutorials: Nested Classes + Exceptions + Concurrency + Basic I/O (accessed \today).; Java SE API docs (e.g., \texttt{Thread}, \texttt{AutoCloseable}) (accessed \today).}

\newcommand{\OOPSUnitIVSources}{Schildt, \textit{Java: The Complete Reference} (12e) (generics, lambdas, Swing, JDBC).; Bloch, \textit{Effective Java} (3e) (generics + lambdas).; Oracle Java Tutorials: Generics + Lambda Expressions + Swing + JDBC Basics (accessed \today).; Java SE API docs (e.g., \texttt{java.util.function}, \texttt{javax.swing}, \texttt{java.sql}) (accessed \today).}

\newcommand{\OOPSUnitVSources}{Schildt, \textit{Java: The Complete Reference} (12e) (Collections Framework + wrapper classes).; Oracle Java docs: Collections Framework overview (accessed \today).; Java SE API docs (\texttt{java.util} collections; wrapper classes in \texttt{java.lang}) (accessed \today).}

\newcommand{\OOPSUnitVISources}{Git SCM: \textit{Pro Git} (2e) (version control workflow), accessed \today.; GitHub Docs (repositories, issues, pull requests), accessed \today.; Oracle Java Tutorials: Swing + JDBC (accessed \today).; JUnit 5 User Guide (accessed \today).; Apache Maven Project: Getting Started (accessed \today).; Gradle User Manual: Getting Started (accessed \today).; UML class diagrams: UML 2.x overview/spec (accessed \today).}

\newcommand{\OOPSMiscWhyInterfacesSources}{Schildt, \textit{Java: The Complete Reference} (12e) (interfaces).; Bloch, \textit{Effective Java} (3e) (prefer interfaces where appropriate).; Oracle Java Tutorials: Interfaces (accessed \today).; Java SE API docs (e.g., \texttt{Comparable}, \texttt{Runnable}) (accessed \today).}



\title{Object Oriented Programming (CSEG1044)\\Unit 06 -- Lecture 01 Notes\\Capstone Kickoff + Team Setup + Git/GitHub}
\author{Tofik Ali}
\date{\today}

\begin{document}
\maketitle
\tableofcontents

\section{Lecture Overview}
This is the start of the capstone project.
The goal is not only ``write code''; the goal is to build a small but complete software product with:
\begin{itemize}
  \item clear scope,
  \item clean OOP design,
  \item a working Swing UI,
  \item JDBC persistence,
  \item and a professional Git/GitHub workflow.
\end{itemize}
\notesources{\OOPSUnitVISources}

\section{Syllabus Alignment (Unit 06)}
\textbf{Unit 06:} standalone Java project; source control; collaboration; build and artifacts

\section{Learning Outcomes}
After this lecture, you should be able to:
\begin{itemize}
  \item Select a realistic problem statement with clear scope boundaries.
  \item Create a team repo and organize work using issues and branches.
  \item Follow a simple Git workflow (branching + pull requests).
  \item Produce Deliverable-0: proposal (1 page) + initial backlog (issues).
\end{itemize}

\section{Project theme options (pick one)}
Choose a domain that naturally needs OOP modelling + collections + Swing UI + JDBC persistence:
\begin{itemize}
  \item Library management (books, members, issue/return)
  \item Hostel / room booking
  \item Inventory management (products, stock in/out)
  \item Clinic appointment booking
  \item Campus resource booking (classrooms/labs/equipment) --- recommended
\end{itemize}

\section{Scope boundaries (most important success factor)}
Projects fail mostly due to unclear scope.
Use a simple scope definition:
\begin{itemize}
  \item \textbf{Must-have}: required for minimum working demo.
  \item \textbf{Should-have}: improves product but not required for demo.
  \item \textbf{Could-have}: optional enhancements if time permits.
  \item \textbf{Non-goals}: explicitly out of scope (avoid wasted effort).
\end{itemize}

\textbf{Example (Campus resource booking):}
\begin{itemize}
  \item Must: add/list resources; create booking; prevent overlapping bookings.
  \item Should: user login (simple roles); booking cancel; search/filter.
  \item Could: export report (CSV/PDF); notifications.
  \item Non-goals: full web app; payment gateway; complex authentication.
\end{itemize}

\section{Evaluation mindset (what usually matters)}
Even if the exact rubric differs, these are typical expectations:
\begin{itemize}
  \item \textbf{Correctness}: features work, edge cases handled.
  \item \textbf{Design}: good class structure; separation of UI vs DB vs logic.
  \item \textbf{Database}: schema + CRUD, data persists correctly.
  \item \textbf{UI}: usable screens, validations, meaningful messages.
  \item \textbf{Git hygiene}: commits, branches, PRs, issues, evidence of collaboration.
  \item \textbf{Documentation}: README, setup steps, screenshots (optional), report/proposal.
\end{itemize}

\section{Team setup checklist}
\begin{enumerate}
  \item Create team (3--4 students, as per batch rule).
  \item Create a GitHub repository (one repo per team).
  \item Add collaborators (all team members + instructor if needed).
  \item Add a \texttt{README.md} with:
  \begin{itemize}
    \item project title + short description,
    \item tech stack (Java, Swing, JDBC, DB),
    \item how to run (steps),
    \item team members.
  \end{itemize}
  \item Create labels: \texttt{feature}, \texttt{bug}, \texttt{doc}, \texttt{db}, \texttt{ui}.
\end{enumerate}

\section{Git/GitHub workflow (simple and effective)}
\subsection{Branching convention}
\begin{itemize}
  \item \texttt{main}: always stable / demo-ready.
  \item Feature branches: \texttt{feature/<short-name>} (e.g., \texttt{feature/resource-crud}).
  \item Fix branches: \texttt{fix/<short-name>}.
\end{itemize}

\subsection{Pull request (PR) habit}
For each feature:
\begin{enumerate}
  \item Create issue.
  \item Create branch from \texttt{main}.
  \item Commit frequently with meaningful messages.
  \item Open a PR and request review from teammate.
  \item Merge after review (and after fixing comments).
\end{enumerate}

\subsection{Commit message guideline}
Keep it short and clear:
\begin{itemize}
  \item \texttt{Add ResourceDao insert + findAll}
  \item \texttt{Fix booking overlap validation}
  \item \texttt{Update README setup steps}
\end{itemize}

\section{How to write good GitHub issues (beginner-friendly)}
GitHub issues are your ``work tracker'' and also your evidence of collaboration.
An issue should be small enough to finish in \textbf{one focused session} (or at most 1--2 sessions).

\subsection{Issue template (simple)}
Use this structure (copy/paste into issue description):
\begin{itemize}
  \item \textbf{Problem/Goal:} what needs to be built or fixed?
  \item \textbf{Acceptance criteria:} 3--6 testable bullets (how do we know it is done?)
  \item \textbf{Tasks:} checklist of steps (UI/service/DAO/DB/test/doc)
  \item \textbf{Notes:} screenshots, UI sketch, edge cases (optional)
\end{itemize}

\subsection{Example issue (Campus Resource Manager)}
\textbf{Title:} Prevent overlapping bookings for same resource\par
\textbf{Acceptance criteria (sample):}
\begin{itemize}
  \item If a new booking overlaps an existing booking for the same resource, the system shows an error message and does not save it.
  \item Edge case: booking that ends at the exact start time of another booking is allowed.
  \item Valid booking is saved to DB and appears in booking list view.
\end{itemize}
\textbf{Tasks (sample):}
\begin{itemize}
  \item Add overlap check method in \texttt{BookingService}
  \item Add DAO query to find overlaps (by resource + time range)
  \item Add UI validation message near booking form
  \item Test with 3 cases (overlap, boundary-touch, non-overlap)
\end{itemize}

\section{Suggested repo structure (keep it simple)}
You can organize your repo like this:
\begin{lstlisting}[language=]
repo-root/
  README.md
  docs/
    deliverable-0-proposal.pdf
    deliverable-1-use-cases.pdf
    deliverable-2-class-diagram.pdf
  src/
    com/upes/crm/...
  db/
    schema.sql
\end{lstlisting}
\textbf{Tip:} Keep deliverables in \texttt{docs/} and keep code in \texttt{src/}. This makes demo preparation easier.

\section{Deliverable-0 (proposal + backlog)}
\subsection*{Deliverable-0A: 1-page proposal (template)}
\begin{itemize}
  \item Title + team members (roll numbers)
  \item Problem statement (3--5 lines)
  \item Must/Should/Could features (bullets)
  \item UI screens (list)
  \item Database entities/tables (list)
  \item Milestones (week-wise or lecture-wise plan)
  \item Risks (2--3) + mitigation
  \item GitHub repo link
\end{itemize}

\subsection*{Deliverable-0B: backlog (initial issues)}
Create issues for each must-have item, plus technical tasks.
Example issue list:
\begin{itemize}
  \item Create DB schema (resources, bookings)
  \item Resource CRUD UI + DAO
  \item Booking create + validation (no overlap)
  \item List bookings (table view)
  \item Exception handling + user-friendly messages
  \item README run instructions
\end{itemize}
\notesources{\OOPSUnitVISources}

\section{In-Class Activity (evidence to produce)}
\begin{itemize}
  \item Repo created + collaborators added
  \item README created
  \item At least 8--12 issues created (must-have + setup tasks)
  \item At least 1 PR opened (even if small: README + folder structure)
\end{itemize}

\section{Checkpoint Questions (with sample answers)}
\textbf{Q1.} Why do we define Must/Should/Could/Non-goals?\par
\textbf{Sample answer:} It prevents scope creep, helps teams finish a working minimum product, and makes planning measurable.

\vspace{0.6em}
\textbf{Q2.} Why use branches + PRs instead of committing everything to \texttt{main}?\par
\textbf{Sample answer:} Branches isolate changes and keep \texttt{main} stable. PRs enable review and reduce bugs, and they create evidence of collaboration.

\end{document}
